\chapter{Modèles continus et méthode d'Euler}
\section*{Exercice 13.1 -- Croissance continue}
La solution est $y(t)=50e^{0{,}2t}$, vérifiée par dérivation et condition initiale ; l'unicité découle du facteur intégrant. Sur une unité, le facteur est $e^{0{,}2}$ et le taux $e^{0{,}2}-1\simeq22{,}14\%$. Le coefficient instantané $0{,}2$ ne correspond pas exactement à $20\%$ par unité.

\section*{Exercice 13.2 -- Décroissance et demi-vie}
$80e^{-0{,}4t}=40$ équivaut à $e^{-0{,}4t}=1/2$, donc $-0{,}4t=-\ln2$. Ainsi $t=\ln2/0{,}4$, dans l'unité de temps du modèle. La stricte décroissance assure l'unicité.

\section*{Exercice 13.3 -- Apport constant}
L'équilibre vaut $3$. Poser $z=y-3$ donne $z'=-2z$ et $z(0)=-2$, donc $y(t)=3-2e^{-2t}$. Elle croît vers $3$ pour $t\ge0$. Vérifier $y'=-2y+6$ contrôle les signes.

\section*{Exercice 13.4 -- Deux pas d'Euler}
Le pas vaut $h=1/2$. Euler donne $Y_1=1+1/2=3/2$, puis $Y_2=(3/2)^2=9/4=2{,}25$. Comme $1+h<e^h$ pour $h>0$, $(1+h)^2<e$, donc l'approximation est par défaut.

\section*{Exercice 13.5 -- Un pas inadéquat}
La récurrence est $Y_{n+1}=(1-3)Y_n=-2Y_n$, donc $Y_n=(-2)^n$. Elle alterne et diverge en valeur absolue. La solution exacte $e^{-3t}$ reste positive et tend vers $0$. Le pas est trop grand pour la stabilité de cette méthode.

\section*{Exercice 13.6 -- Unités du modèle}
$y'$ a l'unité litre par minute. Donc $a$ est en minute$^{-1}$ et $b$ en litre par minute. Avec $s=60t$, la dérivée par rapport à $s$ vaut $(ay+b)/60$ : le nouveau coefficient est $a/60$ en seconde$^{-1}$, et l'apport est $b/60$.
